\begin{table}[H]
\centering
\caption{PIB, ocupados, horas y productividad laboral, 2010--2025}
\label{tab:ocupados_horas_resumen}
\small
\begin{tabular}{lrrr}
\toprule
Indicador & 2010 & 2025 & Crec. anualizado \\
\midrule
PIB real (Billones de pesos de 2015) & 640,2 & 1.021,1 & 3,16\% \\
Ocupados (Millones de personas) & 16,8 & 23,0 & 2,09\% \\
Horas semanales por trabajador (Horas por semana) & 46,8 & 43,7 & -0,45\% \\
PIB por trabajador (Millones de pesos de 2015 por ocupado) & 38,0 & 44,5 & 1,05\% \\
PIB por hora trabajada (Miles de pesos de 2015 por hora) & 15,6 & 19,6 & 1,51\% \\
\bottomrule
\end{tabular}
\caption*{\footnotesize Nota: PIB real en pesos constantes de 2015. Ocupados expandidos con el factor \texttt{fex}. Las horas semanales promedio se calculan como horas anuales totales divididas por ocupados y por 52 semanas. Se excluye 2020 por no contar con GEIH anual comparable en la base del proyecto. Fuente: cálculos propios con DANE, PIB trimestral por el enfoque de producción, y GEIH.}
\end{table}